\sscp{Other classificatinos}{07-01-03}
\citet{Dyen1978b} regrouped Stresemann’s Sub-\ili{Ambon} closer with
Sub-Seram and farther from Sub-\ili{Buru} based mostly on lexicostatistics.
Both \citeauthor{Collins1983}’ (\citeyear{Collins1983}) Nunusaku
and our \tsc{\ili{Ambon}-Seram} subgroup support the conclusion of
Dyen’s adjustment, although unlike Collins we do not link \tsc{\ili{Ambon}-Seram} languages
with \tsc{Sula-Buru} at a higher level.
\citet[45]{Dyen1965} would “perhaps” include \ili{Soboyo} (\tsc{Taliabu})
with his Moluccan linkage.

\citet{Blust1981}, focusing on \ili{Soboyo} (\ili{Taliabu} Island),
grouped the languages of the Sula archipelago in the northwest with \ili{Buru}.
After careful examination, the strongest basis for such a link
is that they share *R > \it{h}.
Elsewhere \citet[279]{Blust1993} suggests that \ili{Geser}
and \ili{Watubela} in far-eastern Seram
and offshore islands should be a “primary branch” of CMP.

In the same work \citet[276ff]{Blust1993} defined a “\ili{Yamdena}-North
Bomberai” subgroup based on a bundle of lexical innovations,
most of which are not exclusive to the group of languages on
which Blust was focused.\footnote{
For example,
reflexes of *bat-bata ‘female’ are found in some \tsc{Seram-Tanimbar-Bomberai},
\tsc{Sula-Buru}, and \tsc{Timor-Babar} languages; *ndundum ‘thunder’
is found in some \tsc{Seram-Tanimbar-Bomberai},
\tsc{Aru}, \tsc{Timor-Babar}, and \tsc{\ili{Bima}-Lembata} languages;
*ŋarak (which we reconstruct as **ŋaRək) ‘year’ is found in some
\tsc{Sula-Buru}, \tsc{\ili{Ambon}-Seram}, \tsc{Aru}, and \tsc{Seram-Tanimbar-Bomberai}
languages; and so forth.}
Yet words like **kəRi ‘tongue’,
while found beyond Blust’s “\ili{Yamdena}-North Bomberai” languages,
nevertheless do seem to be restricted to the wider STB  subgroup.
While the details of Blust’s subgroup itself is problematic,
linking \ili{Yamdena} to the Austronesian languages on the Bomberai
peninsula remains an important contribution.

\citet{Mills1991} argues for a “Tanimbar-\ili{Kei}” subgroup that includes
\ili{Yamdena}, \ili{Fordata}, and \ili{Kei}.
He further linked these languages with \ili{Selaru}
and \ili{Leti} within a “South\ili{east} Maluku” subgroup,
but presented very little data to support it.
His appeal to “binding” (1991:250ff) to further link south\ili{east}
Maluku to \ili{Selaru}, the languages of southwest Maluku,
Timor and west Sumba are not persuasive
(see \citeauthor{Blust1993} [\citeyear[276]{Blust1993}], and \srf{13-02-03}).
\citet{Mills1991} did not investigate the possibility of linking
\ili{Kei}, \ili{Fordata} and \ili{Yamdena} with languages to the north.

\citet{GrayGreenhill2009} include the languages of our \tsc{Taliabu},
\tsc{Sula-Buru}, \tsc{\ili{Ambon}-Seram}, \tsc{Aru} subgroups, and
part of our \tsc{Seram-Tanimbar-Bomberai} subgroup under a single
subgroup they call “\ili{Central} Maluku” (p.480).
Curiously their algorithm places \ili{Buru}
and \ili{Soboyo} as even more divergent from the other languages of
their \ili{Central} Maluku than are the \tsc{Aru} languages (\crf{ch:Aru}),
whereas \citet{Collins1981,Collins1983} repeatedly claims \ili{Buru} is “closely
related” to \ili{Ambon} and Seram languages, and by association,
so is \ili{Soboyo} and \ili{Taliabu}
(\citeyear[15, 19--20]{Collins1983}, \citeyear[86, footnote 1]{Collins1989}).
These two views are not compatible
and raise many questions.
We believe our classification based primarily on the historical
phonology provides a solution. \citet{GrayGreenhill2009}
also group their “\ili{Yamdena}-North Bomberai” separately from
their \ili{Central} Maluku subgroup.
We place those languages within our \tsc{Seram-Tanimbar-Bomberai} subgroup.

\citet{Chlenov1976}, working mostly with lexicostatistics,
included the languages of East Seram,
but not the languages of the Sula archipelago,
within his “\ili{Ambon}” (roughly Collins’ \ili{Central} Maluku) subgroup.
He also placed \ili{Banda} as its own node within his \ili{Ambon} subgroup.
Diverging from several others,
Chlenov placed \ili{Geser} and \ili{Watubela} as their own node under his
separate “South Moluccan” subgroup,
linking them southward to languages such as \ili{Kei}, \ili{Fordata} and \ili{Yamdena}.
This proposal is in alignment with our \tsc{Seram-Tanimbar-Bomberai}.

\citet{Taguchi1989} tabulates the lexical similarity for languages
in west Seram, central Seram, parts of \ili{east} Seram,
and in the west,
several languages in the \ili{Manipa} Strait.
His finding of Boano as being a bit divergent from the rest parallels
our findings in the historical phonology.

\citet{LoskiLoski1989} also working mostly with lexicostatistics,
tabulate the lexical similarity for languages from central Seram,
\ili{east} Seram, and offshore islands including \ili{Geser}-\ili{Gorom},
\ili{Watubela}, \tsc{\ili{Teor}-Kur}, and \ili{Kei}.
Their findings \citeyear[112]{LoskiLoski1989} of \ili{Banda} (Elat) being equally
unrelated to everything else mirror our placement of \ili{Banda},
based primarily on the historical phonology,
as an isolate in the region.\footnote{\label{fn:BandaLexicalSimilarity}
		While \citet{Collins1986}
		groups \ili{Banda} together with \ili{Geser}
		and \ili{Watubela} under his Proto-\ili{Banda},
		\citet[112]{LoskiLoski1989} find \ili{Banda} -- \ili{Geser}-\ili{Gorom} =
		39{\%}, \ili{Banda} -- \ili{Watubela} = 39{\%},
		\ili{Banda} with \tsc{East Seram} languages (\ili{Masiwang},
		\ili{Bobot}) = 35--36{\%}, \ili{Banda} -- \tsc{\ili{Teor}-Kur} = 40{\%},
		\ili{Banda} -- \ili{Kei} = 38{\%}.}
(See \crf{ch:Ban}.)

Focusing further south, \citet{Hughes1987} lexicostatistics study
was restricted geographically in scope to \ili{Kei}, Tanimbar and \ili{Aru}.
He did not investigate how some of those languages might be related
to the languages north of \ili{Kei},
extending to eastern Seram,
a very remote region from any direction.